\documentclass[11pt,a4paper]{article}
\usepackage[utf8]{inputenc}
\usepackage[margin=1in]{geometry}
\usepackage{xcolor}
\usepackage{multicol}
\usepackage{enumitem}
\usepackage{graphicx}

% Define colors matching the text
\definecolor{textpink}{HTML}{E0115F}

\begin{document}

\begin{flushleft}
    {\huge\bfseries\itshape\color{textpink} 1. The Portrait of a Lady} \\[0.5em]
    \hfill {\Large\itshape Khushwant Singh}
\end{flushleft}

\vspace{1.5em}

\noindent \textbf{\sffamily Notice these expressions in the text.} \\
\noindent \textbf{\sffamily Infer their meaning from the context.}

\vspace{1em}

\begin{multicols}{2}
\begin{itemize}[label=\(\bullet\), itemsep=0.3em]
    \item the thought was almost revolting
    \item an expanse of pure white serenity
    \item[$\diamond$] a turning-point
    \item[$\diamond$] accepted her seclusion with resignation
    \item[\(\bullet\)] a veritable bedlam of chirpings
    \item[\(\bullet\)] frivolous rebukes
    \item[\(\bullet\)] the sagging skins of the dilapidated drum
\end{itemize}
\end{multicols}

\vspace{2em}

My grandmother, like everybody's grandmother, was an old woman. She had been old and wrinkled for the twenty years that I had known her. People said that she had once been young and pretty and had even had a husband, but that was hard to believe. My grandfather's portrait hung above the mantelpiece in the drawing room. He wore a big turban and loose-fitting clothes. His long, white beard covered the best part of his chest and he looked at least a hundred years old. He did not look the sort of person who would have a wife or children. He looked as if he could only have lots and lots of grandchildren. As for my grandmother being young and pretty, \textbf{the thought was almost revolting.} She often told us of the games she used to play as a child. That seemed quite absurd and undignified on her part and we treated it like the fables of the Prophets she used to tell us.

She had always been short and fat and slightly bent. Her face was a criss-cross of wrinkles running from everywhere to everywhere. No, we were certain she had always been as we had

\vfill
\begin{center}
    \small 2024-25
\end{center}

\end{document}
\documentclass[11pt,a4paper]{article}
\usepackage[utf8]{inputenc}
\usepackage[margin=1in]{geometry}
\usepackage{xcolor}
\usepackage{microtype}

\begin{document}

% Page Header
\noindent
4 \hfill \textsc{Hornbill}

\vspace{1.5em}

\begin{flushleft}
    \noindent known her. Old, so terribly old that she could not have grown older, and had stayed at the same age for twenty years. She could never have been pretty; but she was always beautiful. She hobbled about the house in spotless white with one hand resting on her waist to balance her stoop and the other telling the beads of her rosary. Her silver locks were scattered untidily over her pale, puckered face, and her lips constantly moved in inaudible prayer. Yes, she was beautiful. She was like the winter landscape in the mountains, \textbf{an expanse of pure white serenity} breathing peace and contentment.
    
    \vspace{1em}
    
    My grandmother and I were good friends. My parents left me with her when they went to live in the city and we were constantly together. She used to wake me up in the morning and get me ready for school. She said her morning prayer in a monotonous sing-song while she bathed and dressed me in the hope that I would listen and get to know it by heart; I listened because I loved her voice but never bothered to learn it. Then she would fetch my wooden slate which she had already washed and plastered with yellow chalk, a tiny earthen ink-pot and a red pen, tie them all in a bundle and hand it to me. After a breakfast of a thick, stale chapatti with a little butter and sugar spread on it, we went to school. She carried several stale chapattis with her for the village dogs.
    
    \vspace{1em}
    
    My grandmother always went to school with me because the school was attached to the temple. The priest taught us the alphabet and the morning prayer. While the children sat in rows on either side of the verandah singing the alphabet or the prayer in a chorus, my grandmother sat inside reading the scriptures. When we had both finished, we would walk back together. This time the village dogs would meet us at the temple door. They followed us to our home growling and fighting with each other for the chapattis we threw to them.
    
    \vspace{1em}
    
    When my parents were comfortably settled in the city, they sent for us. That was \textbf{a turning-point} in our friendship. Although we shared the same room, my grandmother no longer came to school with me. I used to go to an English school in a motor bus. There were no dogs in the streets and she took to feeding sparrows in the courtyard of our city house.
    
    \vspace{1em}
    
    As the years rolled by we saw less of each other. For some time she continued to wake me up and get me ready for school. When I came back she would ask me what the teacher had
\end{flushleft}

\end{document}
\documentclass[11pt,a4paper]{article}
\usepackage[utf8]{inputenc}
\usepackage[margin=1in]{geometry}
\usepackage{xcolor}
\usepackage{microtype}

\begin{document}

% Page Header
\noindent
\textsc{The Portrait of a Lady} \hfill 5

\vspace{1.5em}

\begin{flushleft}
    \noindent taught me. I would tell her English words and little things of western science and learning, the law of gravity, Archimedes' Principle, the world being round, etc. This made her unhappy. She could not help me with my lessons. She did not believe in the things they taught at the English school and was distressed that there was no teaching about God and the scriptures. One day I announced that we were being given music lessons. She was very disturbed. To her music had lewd associations. It was the monopoly of harlots and beggars and not meant for gentlefolk. She said nothing but her silence meant disapproval. She rarely talked to me after that.
    
    \vspace{1em}
    
    When I went up to University, I was given a room of my own. The common link of friendship was snapped. My grandmother \textbf{accepted her seclusion with resignation}. She rarely left her spinning-wheel to talk to anyone. From sunrise to sunset she sat by her wheel spinning and reciting prayers. Only in the afternoon she relaxed for a while to feed the sparrows. While she sat in the verandah breaking the bread into little bits, hundreds of little birds collected round her creating \textbf{a veritable bedlam of chirrupings}. Some came and perched on her legs, others on her shoulders. Some even sat on her head. She smiled but never shooed them away. It used to be the happiest half-hour of the day for her.
    
    \vspace{1em}
    
    When I decided to go abroad for further studies, I was sure my grandmother would be upset. I would be away for five years, and at her age one could never tell. But my grandmother could. She was not even sentimental. She came to leave me at the railway station but did not talk or show any emotion. Her lips moved in prayer, her mind was lost in prayer. Her fingers were busy telling the beads of her rosary. Silently she kissed my forehead, and when I left I cherished the moist imprint as perhaps the last sign of physical contact between us.
    
    \vspace{1em}
    
    But that was not so. After five years I came back home and was met by her at the station. She did not look a day older. She still had no time for words, and while she clasped me in her arms I could hear her reciting her prayers. Even on the first day of my arrival, her happiest moments were with her sparrows whom she fed longer and with \textbf{frivolous rebukes}.
    
    \vspace{1em}
    
    In the evening a change came over her. She did not pray. She collected the women of the neighbourhood, got an old drum and started to sing. For several hours she thumped \textbf{the sagging}
\end{flushleft}

\end{document}
\documentclass[11pt,a4paper]{article}
\usepackage[utf8]{inputenc}
\usepackage[margin=1in]{geometry}
\usepackage{xcolor}
\usepackage{microtype}

% Define colors matching the original document highlights
\definecolor{magenta-accent}{HTML}{E0115F}

\begin{document}

% Page Header
\noindent
6 \hfill \textsc{Hornbill}

\vspace{1.5em}

\begin{flushleft}
    \noindent \textbf{skins of the dilapidated drum} and sang of the home-coming of warriors. We had to persuade her to stop to avoid overstraining. That was the first time since I had known her that she did not pray.
    
    \vspace{1em}
    
    The next morning she was taken ill. It was a mild fever and the doctor told us that it would go. But my grandmother thought differently. She told us that her end was near. She said that, since only a few hours before the close of the last chapter of her life she had omitted to pray, she was not going to waste any more time talking to us.
    
    \vspace{1em}
    
    We protested. But she ignored our protests. She lay peacefully in bed praying and telling her beads. Even before we could suspect, her lips stopped moving and the rosary fell from her lifeless fingers. A peaceful pallor spread on her face and we knew that she was dead.
    
    \vspace{1em}
    
    We lifted her off the bed and, as is customary, laid her on the ground and covered her with a red shroud. After a few hours of mourning we left her alone to make arrangements for her funeral. In the evening we went to her room with a crude stretcher to take her to be cremated. The sun was setting and had lit her room and verandah with a blaze of golden light. We stopped half-way in the courtyard. All over the verandah and in her room right up to where she lay dead and stiff wrapped in the red shroud, thousands of sparrows sat scattered on the floor. There was no chirping. We felt sorry for the birds and my mother fetched some bread for them. She broke it into little crumbs, the way my grandmother used to, and threw it to them. The sparrows took no notice of the bread. When we carried my grandmother's corpse off, they flew away quietly. Next morning the sweeper swept the bread crumbs into the dustbin.
    
    \vspace{1.5em}
    \rule{\linewidth}{0.5pt}
    \vspace{1.5em}

    % Heading Section
    {\Large\color{black!80} Understanding the text \hspace{0.2em}\color{magenta-accent!60}$\blacksquare$}
    
    \vspace{1.5em}
    
    \noindent Mention
    \begin{enumerate}
        \item The three phases of the author's relationship with his grandmother before he left the country to study abroad.
        \item Three reasons why the author's grandmother was disturbed when he started going to the city school.
    \end{enumerate}
\end{flushleft}

\end{document}
\documentclass[11pt,a4paper]{article}
\usepackage[utf8]{inputenc}
\usepackage[margin=1in]{geometry}
\usepackage{xcolor}
\usepackage{amssymb}
\usepackage{microtype}

% Define accent color matching the document's design
\definecolor{magenta-accent}{HTML}{E0115F}

\begin{document}

% Page Header
\noindent
\textsc{The Portrait of a Lady} \hfill 7

\vspace{1.5em}

\begin{flushleft}
    % Continuation of previous list
    \begin{enumerate}
        \setcounter{enumi}{2}
        \item Three ways in which the author's grandmother spent her days after he grew up.
        \item The odd way in which the author's grandmother behaved just before she died.
        \item The way in which the sparrows expressed their sorrow when the author's grandmother died.
    \end{enumerate}
    
    \vspace{2em}

    % Heading 1
    {\Large\color{black!80} Talking about the text \hspace{0.2em}\color{magenta-accent!60}$\blacksquare$}
    
    \vspace{1em}
    \noindent Talk to your partner about the following.
    \begin{enumerate}
        \item The author's grandmother was a religious person. What are the different ways in which we come to know this?
        \item Describe the changing relationship between the author and his grandmother. Did their feelings for each other change?
        \item Would you agree that the author's grandmother was a person strong in character? If yes, give instances that show this.
        \item Have you known someone like the author's grandmother? Do you feel the same sense of loss with regard to someone whom you have loved and lost?
    \end{enumerate}
    
    \vspace{2em}

    % Heading 2
    {\Large\color{black!80} Thinking about language \hspace{0.2em}\color{magenta-accent!60}$\blacksquare$}
    
    \vspace{1.5em}
    \begin{enumerate}
        \item Which language do you think the author and his grandmother used while talking to each other?
        \item Which language do you use to talk to elderly relatives in your family?
        \item How would you say `a dilapidated drum' in your language?
        \item Can you think of a song or a poem in your language that talks of homecoming?
    \end{enumerate}
    
    \vspace{2em}

    % Heading 3
    {\Large\color{black!80} Working with words \hspace{0.2em}\color{magenta-accent!60}$\blacksquare$}
    
    \vspace{1.5em}
    \noindent I. Notice the following uses of the word `tell' in the text.
    \begin{enumerate}
        \item Her fingers were busy \textit{telling the beads} of her rosary.
        \item I would \textit{tell her} English words and little things of Western science and learning.
        \item At her age \textit{one could never tell}.
        \item She \textit{told us} that her end was near.
    \end{enumerate}
\end{flushleft}

\end{document}
\documentclass[11pt,a4paper]{article}
\usepackage[utf8]{inputenc}
\usepackage[margin=1in]{geometry}
\usepackage{xcolor}
\usepackage{amssymb}
\usepackage{array}
\usepackage{microtype}

% Define accent color matching the document's design
\definecolor{magenta-accent}{HTML}{E0115F}

\begin{document}

% Page Header
\noindent
8 \hfill \textsc{Hornbill}

\vspace{1.5em}

\begin{flushleft}
    \noindent Given below are four different senses of the word `tell'. Match the meanings to the uses listed above.
    \begin{enumerate}
        \item make something known to someone in spoken or written words
        \item count while reciting
        \item be sure
        \item give information to somebody
    \end{enumerate}
    
    \vspace{1.5em}
    
    \noindent II. Notice the different senses of the word `take'.
    \begin{enumerate}
        \item to \textit{take to} something: to begin to do something as a habit
        \item to \textit{take ill}: to suddenly become ill
    \end{enumerate}
    \noindent Locate these phrases in the text and notice the way they are used.
    
    \vspace{1.5em}
    
    \noindent III. The word `hobble' means to walk with difficulty because the legs and feet are in bad condition.
    
    \vspace{0.5em}
    \noindent Tick the words in the box below that also refer to a manner of walking.
    
    \vspace{1em}
    \noindent
    \begin{tabular}{|p{0.18\linewidth} p{0.18\linewidth} p{0.18\linewidth} p{0.18\linewidth} p{0.18\linewidth}|}
        \hline
        haggle   & shuffle & stride  & ride   & waddle \\
        wriggle  & paddle  & swagger & trudge & slog   \\
        \hline
    \end{tabular}
    
    \vspace{2.5em}

    % Heading
    {\Large\color{black!80} Noticing form \hspace{0.2em}\color{magenta-accent!60}$\blacksquare$}
    
    \vspace{1.5em}
    
    \noindent Notice the form of the verbs italicised in these sentences.
    \begin{enumerate}
        \item My grandmother was an old woman. She \textit{had been} old and wrinkled for the twenty years that I \textit{had known} her. People said that she \textit{had} once \textit{been} young and pretty and \textit{had} even \textit{had} a husband, but that was hard to believe.
        \item When we both \textit{had finished} we would walk back together.
        \item When I came back she would ask me what the teacher \textit{had taught} me.
        \item It was the first time since I \textit{had known} her that she did not pray.
        \item The sun was setting and \textit{had lit} her room and verandah with a golden light.
    \end{enumerate}
    
    \vspace{1em}
    \noindent These are examples of the past perfect forms of verbs. When we recount things in the distant past we use this form.
\end{flushleft}

\end{document}
\documentclass[11pt,a4paper]{article}
\usepackage[utf8]{inputenc}
\usepackage[margin=1in]{geometry}
\usepackage{xcolor}
\usepackage{amssymb}
\usepackage{microtype}

% Define accent color matching the document's design
\definecolor{magenta-accent}{HTML}{E0115F}

\begin{document}

% Page Header
\noindent
\textsc{The Portrait of a Lady} \hfill 9

\vspace{1.5em}

\begin{flushleft}
    % Section 1
    {\Large\color{black!80} Things to do \hspace{0.2em}\color{magenta-accent!60}$\blacksquare$}
    
    \vspace{1em}
    \noindent Talk with your family members about elderly people who you have been intimately connected with and who are not there with you now. Write a short description of someone you liked a lot.
    
    \vspace{2.5em}
    
    % Notes Divider Heading
    \begin{center}
        {\Large\color{magenta-accent} \rule[0.5ex]{2cm}{1pt} Notes \rule[0.5ex]{2cm}{1pt}}
    \end{center}
    
    \vspace{1.5em}

    % Sub-heading 1
    {\large\color{black!80} Understanding the text \hspace{0.2em}\color{magenta-accent!60}$\blacksquare$}
    
    \vspace{0.8em}
    \noindent The tasks cover the entire text and help in summarising the various phases of the autobiographical account and are based on the facts presented.
    
    \begin{itemize}
        \item Ask the students to read the text silently, paragraph by paragraph, and get a quick oral feedback on what the main points of each are. For example: Para 1 -- description of grandmother and grandfather's photograph.
        \item At the end of the unit ask students to answer the comprehension questions first orally and then in writing in point form.\\
        For example, when he went to the:
        \begin{itemize}
            \item[--] village school
            \item[--] city school
            \item[--] university
        \end{itemize}
    \end{itemize}
    
    \vspace{1.5em}

    % Sub-heading 2
    {\large\color{black!80} Talking about the text \hspace{0.2em}\color{magenta-accent!60}$\blacksquare$}
    
    \vspace{0.8em}
    \noindent Peer interaction about the text is necessary before students engage in writing tasks. The questions raised in this section elicit subjective responses to the facts in the text and also open up possibilities for relating the events to the reader's own life and establish the universality of the kind of relationship and feelings described in the text.
    
    \vspace{1.5em}

    % Sub-heading 3
    {\large\color{black!80} Thinking about language \hspace{0.2em}\color{magenta-accent!60}$\blacksquare$}
    
    \vspace{0.8em}
    \noindent The questions here try to:
    \begin{itemize}
        \item make the reader visualise the language that must have been used by the author and his grandmother
        \item think about their own home language
    \end{itemize}
\end{flushleft}

\end{document}
\documentclass[11pt,a4paper]{article}
\usepackage[utf8]{inputenc}
\usepackage[margin=1in]{geometry}
\usepackage{xcolor}
\usepackage{amssymb}
\usepackage{microtype}

% Define accent color matching the document's design
\definecolor{magenta-accent}{HTML}{E0115F}

\begin{document}

% Page Header
\noindent
10 \hfill \textsc{Hornbill}

\vspace{1.5em}

\begin{flushleft}
    % Continuation of previous itemized list
    \begin{itemize}
        \item find equivalents in their language for English phrases
        \item relate to songs with emotional import in their own language.
    \end{itemize}
    
    \vspace{2em}

    % Section 1
    {\large\color{black!80} Working with words \hspace{0.2em}\color{magenta-accent!60}$\blacksquare$}
    
    \vspace{0.8em}
    \noindent Highlight different uses of common words like `tell' and `take'; words used for different ways of walking; and semantically-related word groups. You could add to the items by using the dictionary for vocabulary enrichment.
    
    \vspace{1.5em}

    % Section 2
    {\large\color{black!80} Noticing form \hspace{0.2em}\color{magenta-accent!60}$\blacksquare$}
    
    \vspace{0.8em}
    \noindent Make students notice the use of the past perfect form of the verb that frequently appear in the text to recount the remote past. You could practise the form with other examples.
    
    \vspace{1.5em}

    % Section 3
    {\large\color{black!80} Things to do \hspace{0.2em}\color{magenta-accent!60}$\blacksquare$}
    
    \vspace{0.8em}
    \noindent Relating the topic of the text to the reader's real-life experience; writing about a person who one holds dear.
\end{flushleft}

\vfill

% Page Footer
\begin{center}
    \small 2024-25
\end{center}

\end{document}

